\section{Exact isomorphisms of geometric systems}
\label{exact:section-boundaries}

We continue with the pointed pregeometry of
Theorem~\ref{geom:pregeometry}, above the threshold of
Corollary~\ref{geom:tail-no-growth}.  A shape and its active labels have the
meaning of Definition~\ref{sys:shapes}.  Unless otherwise stated, every
active basis piece in this section has cardinality equal to the uniform
level~$\mu$.  The bases remain auxiliary choices.

\begin{definition}[A prescribed boundary isomorphism]
\label{exact:boundary-data}
Let $a$ be a finite ordered set of axes, and let
$\mathcal X^-=(X_s^-:s\subseteq a)$ and
$\mathcal X^+=(X_s^+:s\subseteq a)$ be systems of the same shape.
A prescribed boundary isomorphism is a family
\[
 h_s:X_s^-\longrightarrow X_s^+\qquad(s\subsetneq a)
\]
of pointed isomorphisms such that
$h_t\mathbin{\upharpoonright}X_s^-=h_s$ whenever
$s\subseteq t\subsetneq a$.  This condition refers to the displayed
inclusions.  It does not add an assumption about intersections of two
incomparable vertices inside the top.
\end{definition}

A list on the boundary can consequently be specified by assigning a
proper face to each of its occurrences.  Before consistency on actual
intersections has been proved, its prescribed image means the list
obtained by applying the corresponding $h_s$ to each occurrence.
In particular, the assertion that two occurrences name the same element
is one of the atomic comparisons that must be preserved.

\begin{theorem}[Exact geometric boundary lifting]
\label{exact:geometric-boundary}
For every finite allowed shape, at all sufficiently closed uniform
levels, every prescribed boundary isomorphism between two systems of
that shape extends to a pointed isomorphism of their tops.  The top
isomorphism, together with its prescribed restrictions, is an
isomorphism of the systems.

The refinements of the eligible levels needed for a fixed shape use
only the one-dimensional bounds for its working scheme, the bounds
for its tiling layouts, and the corresponding assertions for shapes
on fewer axes.  In particular these refinements have the uniformity
specified in Section~\ref{close:section-uniform}.
\end{theorem}

We prove the theorem by induction on the number of axes.  At zero axes
it is pointed categoricity.  The induction at a positive number of
axes has two parts: preserving the joint diagrams of the entire
boundary, and realizing successive extensions of a partial match
inside the given target.  The first part supplies, among other
things, the missing compatibility on actual intersections.

\begin{lemma}[Joint boundary comparisons]
\label{exact:joint-boundary}
Assume exact geometric boundary lifting for all allowed shapes on
fewer than $n$ axes.  At a sufficiently closed uniform level, a
prescribed boundary isomorphism between two $n$-axis systems preserves
the bounded $\K$-diagram of every occurrence-tagged boundary list,
computed in the respective top models.
\end{lemma}

\begin{proof}
If the top label is inactive, let $u$ be the largest active label.
It is a proper subset of the axis set, and the top vertex is $X_u$.
Every vertex is the vertex at its largest active lower label, which
is contained in $u$.  Thus the map at $u$ already gives all the
comparisons in question.  We may assume the top label is active.

Write the axis set as $a\cup\{k\}$, with $k$ last and $|a|=n-1$.
Splitting off $k$ writes the two systems as
\[
 X^-\longrightarrow Y^- ,\qquad X^+\longrightarrow Y^+
\]
in the scheme $\mathcal G_L$ on $a$.
Its categoricity is supplied by
Theorem~\ref{sys:basis-classification}; after refining the levels,
Proposition~\ref{sys:working-package} supplies its working package.
The arrows add rank $\mu$ at the labels in a fixed upward closed
$Q\subseteq L$ and add no rank at the other labels.  Call such a
step a step of shape~$Q$.  The boundary prescription gives a full
system isomorphism $e:X^-\simeq X^+$ and compatible isomorphisms
\[
 h_s:Y_s^-\simeq Y_s^+\qquad(s\subsetneq a)
\]
whose restrictions to $X_s^-$ agree with~$e$.

If joint comparison fails, choose bounded lists $f$ from $X^-$ and
$d^-$ from the proper vertices of $Y^-$ witnessing failure, and put
$d^+=h(d^-)$.  All occurrences from the lower layer may be included
in $f$; all the remaining occurrences lie in these upper proper
vertices.  The comparisons of
\begin{equation}
 (f,d^-)\text{ in }Y^-
 \quad\text{and}\quad
 (e(f),d^+)\text{ in }Y^+
 \label{exact:opposite-comparisons}
\end{equation}
differ.  The notation in \eqref{exact:opposite-comparisons} includes
the indicated vertex occurrences.  If there is no upper test $d^-$,
the base isomorphism and the strong arrows already give equality,
so this case cannot witness failure.

A system isomorphism matching these lists would in particular be a
$\K$-isomorphism of the tops matching them.  Hence the working
$\mathcal G_L$-diagrams of the lists in
\eqref{exact:opposite-comparisons} also differ.  Choose an infinite
cardinal $\sigma<\mu$ bounding the lengths of the lists.  Enlarge a
cardinal $\chi<\mu$ so that $\chi^\sigma=\chi$, it bounds the
relevant labels and parameter-orbit constants, so that
Lemma~\ref{calc:parameter-orbits} gives at most $\chi$ orbits of
$\sigma$-lists over $\chi$ parameters.  Choose $\mu$ closed far enough that all
the following tree costs are below it.  Let $\epsilon$ be the least
infinite cardinal with $2^\epsilon>\chi$.  There are at most $\chi$
nodes below height $\epsilon$, and $2^\epsilon<\mu$ leaves.

We construct a binary tree of branch worlds $C_\eta$ in
$\mathcal G_L$, for $\eta\in{}^{\le\epsilon}2$, and a single
increasing reference chain $(R_i:i\le\epsilon)$ in that scheme.
All these objects have level~$\mu$.  A branch world carries a copy
$f_\eta$ of~$f$.  For every node $\eta$ of length $i$ the invariant
also includes a prescribed isomorphism
\begin{equation}
 b_\eta:\partial C_\eta\simeq\partial R_i.
 \label{exact:reference-invariant}
\end{equation}
Along a branch, these boundary maps commute with the inclusions into
later branch and reference worlds.  At each splitting node, the two
new copies of its upper test $d$ have the same image in the next
reference boundary.  Those images will be retained at all later
stages.

Start with $C_\varnothing=R_0=X^-$ and the identity boundary map.
Suppose the objects at level $i$ have been constructed.  Choose a
reference step $R_i\to R_{i+1}$ of shape~$Q$.  Such a step exists:
categoricity of $\mathcal G_L$ identifies $R_i$ with $X^-$, and we
transport the given extension $X^-\to Y^-$ along that
identification.

For each $\eta\in{}^i2$, an isomorphism $X^-\simeq C_\eta$ can
be chosen to take $f$ to $f_\eta$.  Indeed arrows along the branch
preserve its bounded diagram, and equality of bounded diagrams in
the working scheme is implemented by an isomorphism.  Transport
the minus and plus extensions along this return isomorphism, using
$e$ on the plus base.  They give children $C_{\eta^\frown 0}$ and
$C_{\eta^\frown 1}$ over $C_\eta$, with opposite comparisons to
$f_\eta$ and with the given identification between their proper
upper vertices.  On their lower proper vertices that identification
is the identity of the current branch world.

It remains to extend \eqref{exact:reference-invariant} to the
children; merely choosing arbitrary isomorphisms of the new worlds
would not suffice.  First do this for the plus child.  Process all
proper subsets $s\subsetneq a$ in increasing cardinality.  Consider
the restriction of the plus step to the axes $s\cup\{k\}$ and the
corresponding restriction of $R_i\to R_{i+1}$.  On the base
vertex, the prescribed map is $(b_\eta)_s$.  On every other proper
vertex it is either a restriction of that map or one of the upper
maps already constructed at a proper subset of~$s$.  These maps
commute on all indicated subfaces by construction.  The restriction
has at most $n-1$ axes, so the induction hypothesis lifts this
boundary prescription onto its upper vertex.  This supplies the
map at $s$, whether that vertex is active or repeated from a
smaller active one.  The resulting family is the required plus
boundary map.  Compose it with the given minus-to-plus boundary
identification to obtain the minus boundary map.  Both extend
$b_\eta$, and the two copies of the test have equal reference
images.

At a limit $i\le\epsilon$, take unions along each branch and in
the reference chain.  The lower-dimensional working scheme has
the chain and common-upper-bound properties, and these unions have
size $\mu$.  The coherent boundary maps have unions which are
onto the reference boundary: each previous reference vertex was
already the image of the corresponding branch vertex.  Thus the
invariant persists.  No assertion that a limit inclusion is itself
a step of shape~$Q$ is needed.  At its next successor stage a fresh
step of that shape is available over the union by categoricity.

At a leaf, apply the induction hypothesis once more, this time to
the $(n-1)$-axis worlds $C_\eta$ and $R_\epsilon$, with the
boundary prescription $b_\eta$.  Obtain a full isomorphism
$j_\eta:C_\eta\simeq R_\epsilon$.  Let $A$ be the list of all
reference images of the splitting tests.  Its length is at most
$\chi$, since there were at most $\chi$ splitting nodes and each
test had length at most~$\sigma$.

For two different leaves, use their first splitting node.  The
associated test has the same image in $A$ on both branches, whereas
its comparison with the images of $f$ is opposite.  Arrow
preservation of diagrams carries this distinction to the leaves,
and $j_\eta$ preserves it in $R_\epsilon$.  Consequently the
$2^\epsilon$ lists $j_\eta(f_\eta)$ lie in different orbits over
$A$.  This contradicts the bound~$\chi$.  Joint boundary
comparison follows.
\end{proof}

\begin{remark}
The proof just given applies to equalities between occurrences from
incomparable faces.  It therefore proves that the prescribed maps
have a well-defined injective union on the actual union of the
boundary vertices.  Applying the same observation to the inverse
family proves that this union is a bijection between the two
boundary unions.  Equality of actual intersections was not an
input to the argument.
\end{remark}

The boundary prescription now respects actual intersections and all
bounded joint diagrams.  The remaining task is to extend a bounded match
inside the prescribed target.  We first establish testing for its tiled
boundary layout, then use a strip avoided by all boundary supports to
realize the resulting comparison without enlarging that target.

\begin{lemma}[Testing on rationally tiled parameter orders]
\label{exact:tiled-testing}
Fix the finite-support order template $E$ and its natural stabilized
$\K$-diagrams.  Fix a nonempty set-sized initial order $J_0$, finitely many tile
labels $i<m$, and a finite family $\mathcal A$ of subsets of $m$.
These are fixed descriptor data.  For well-orders $H_i$, put
\[
 J=J_0+\sum_{i<m}(H_i\times\mathbb Q),
\]
with lexicographic order on each product.  For $A\in\mathcal A$, let
$J_A$ consist of $J_0$ and the tiles whose labels belong to $A$.
Let $B$ be the union of the designated term sets $E(J_A)$, with the
term representations inherited from $E$.  Occurrences may retain
their designated-part tags.

There is a monotone cardinal bound $g_{\mathcal A}(\sigma,\tau)$,
depending also on $E$ and $J_0$, with the following property whenever
$|J|=\mu$ and $\mu$ is a sufficiently closed level above the input
cardinals, $|J_0|$, and this bound.  Let $c$ be a
fixed tuple in $E(J)$ of length at most $\tau$, with fixed term
representations.  Two placements of tuples of length at most
$\sigma$ in arbitrary order enlargements of $J$ have the same
comparison over $B$ together with the fixed tuple $c$ if they agree
on every parameter query from $B$ of size at most
$g_{\mathcal A}(\sigma,\tau)$, with $c$ retained as fixed data.
Here $\sigma$ and $\tau$ may be increased to infinite cardinals.
The bound is independent of the horizontal order lengths and of
$\mu$, and is uniform over a fixed set of descriptors of bounded
hereditary size.  The candidate placements need not be terms of a
horizontal reindexing of the tiled template.
\end{lemma}

\begin{proof}
Amalgamate the two supporting order enlargements over $J$ and discard
all added positions not supporting either candidate.  Let $T$ be the
union of the supporting positions of the two candidates and the
chosen old supports of $c$.  Put
\[
 \zeta=\max\bigl(\sigma,\tau,|J_0|,|E|,\aleph_0\bigr).
\]
Thus $|T|\leq\zeta$.  A fiber means the old convex copy
$\{h\}\times\mathbb Q$ at a horizontal position in some $H_i$.
Call an old fiber exceptional if it contains a position of $T$, or
if some added position of $T\setminus J$ has old fiber positions
strictly below and strictly above it.  Each position of $T$ hits or
splits at most one old fiber: two distinct fibers are disjoint convex
suborders.  Hence at most $\zeta$ fibers are exceptional.

Freeze the entire initial order $J_0$, every rational position of
every exceptional fiber, and all positions of $T$.  The resulting
ordered set $F$ has cardinality at most $\zeta$, since every fiber
is countable.  Also freeze the following record: which positions of
$F$ are old, the tile label and rational coordinate of each old
position in an exceptional fiber, the exceptional-fiber grouping,
the labels and supporting lists of both candidates and $c$, and all
order comparisons and equalities among these positions.  The fixed
copy of $J_0$ is retained literally.  Up to renaming the other
positions on a set of size at most $\zeta$, these records form a
set.  In particular, arbitrary cuts made by added supports within a
rational fiber are part of this bounded record; those supports have
not been replaced by horizontal-template terms.

Every nonexceptional old fiber lies wholly in one strict cut of
$F$: it contains no point of $F$, no candidate support splits it,
and distinct exceptional fibers are convex and disjoint from it.
For each tile label $i$ and strict cut $u$ of $F$, record the ordinal
length $\alpha_{i,u}$ of the horizontal positions of
nonexceptional fibers of tile $i$ lying in that cut.  These positions
form an interval of the well-order $H_i$.  There is a set of
coordinates $(i,u)$; the set of cuts can have cardinality as large
as $2^{|F|}$, which still has a bound depending only on $\zeta$.
We do not assume that the supporting order $T$, or the set of its
cuts, is discrete or well-ordered.

For a fixed frozen record, retain only ordinal vectors giving valid
configurations of the displayed tiled form.  The vector and frozen
record determine the parameter order and both candidate support
patterns up to the specified order identifications.  If two valid
vectors satisfy $\alpha_{i,u}\leq\beta_{i,u}$ at every coordinate,
embed the shorter ordinal interval into the longer one at each
coordinate, and send each full rational fiber by the identity on
its rational coordinate.  Together with the identity on $F$, these
maps give an order embedding of the two configurations, fixing the
candidate supports and the old frozen data.  Every old position
stays in its tile, and $J_0$ stays fixed.  Consequently a designated
term remains a term of the same designated part $E(J_A)$, with the
same occurrence tag.  Naturality of the stabilized $\K$-diagrams
preserves each joint diagram comparison under this embedding.

Consider the definable class of valid vectors for which the two
candidate comparisons differ on some parameter query.  Diagrams on
any particular set-sized query can be evaluated in a sufficiently
large eligible padded presentation, as in
Proposition~\ref{lev:testing}.  Thus this is a first-order definable
class with the fixed record as a set parameter.  Apply
Lemma~\ref{lev:ordinal-vectors} to obtain a set of valid unequal
vectors such that every valid unequal vector dominates one of
them.  No assertion that an arbitrary ordinal vector describes a
valid horizontal well-order is needed.  Coordinatewise domination
between the two valid configurations gives the embedding just
constructed, so the query witnessing inequality in the smaller
configuration still witnesses inequality after that embedding.

The old parameter orders in this set of witnesses have a common
cardinal bound.  The number of designated finite-support terms on
one such order is at most its cardinality plus $|E|+\aleph_0$.
Bound this cardinal over the set of witnesses and then over the set
of frozen records.  An unequal comparison in any valid
configuration therefore differs on a designated parameter query of
at most that bounded cardinality, retaining $c$ as fixed data.
Increase the bound to dominate the input cardinals and take
monotone majorants in $\sigma$ and $\tau$.  This defines the required
$g_{\mathcal A}$ independently of the level.  Taking the same
supremum over a fixed set of bounded descriptors proves the last
uniformity assertion.
\end{proof}

We next record the realization step, with its support calculation
explicit.  A permitted parameter layout means either a natural
layout covered by Proposition~\ref{calc:testing} or a rationally
tiled designated-part layout covered by
Lemma~\ref{exact:tiled-testing}.  A bounded list of additional fixed
parameters is allowed in either case.  A testing bound for a
permitted layout will always mean the bound supplied by the
corresponding one of these two results.

\begin{lemma}[Realization using an unoccupied strip]
\label{exact:internal-realization}
Let $E(J)$ be a pointed $\K$-presentation of cardinality $\mu$ at
a sufficiently closed uniform level.  Let $V\subseteq J$ be a
convex substrip of cardinality~$\mu$.  Let $B\subseteq E(J)$ be
a permitted parameter layout, with chosen representations of all its
entries using positions in $J\setminus V$.  Include the fixed
point data among these represented parameters.

Suppose a bounded tuple $z^*$ in a presentation $E(J^*)$, where
$J\subseteq J^*$, has a placed comparison over $B$.  Then a tuple
$z$ of $E(J)$ has the same comparison over every bounded list
from~$B$.  Any specified bounded old entries of $z^*$ are retained
as actual entries of $z$, provided their chosen representations
avoid~$V$.
\end{lemma}

\begin{proof}
Shrink $J^*$ to $J$ together with supports for $z^*$, so its
cardinality is still~$\mu$.  Include the old entries to be retained
in the tuple under consideration and in the bounded fixed data of
the layout.  Choose a bounded test cardinal $\chi<\mu$ at least
the tuple length and the testing bound for this permitted layout.
Let $T^*$ be the chosen supporting positions of $z^*$ in $J^*$.
For the invocation of Lemma~\ref{calc:support-closure}, replace $V$
by its convex hull in $J^*$.  Since $V$ is convex in $J$, this hull
meets $J$ in exactly $V$.  Thus its old complement is still
$J\setminus V$, and the side of this convex hull on the old
complement is the side of the original strip.  The support-closure
lemma therefore has the required convex strip in the ambient
order; the actual realization still uses $E(S\cup V)$.  By
Lemma~\ref{calc:support-closure}, there is a bounded
$S\subseteq J\setminus V$ containing the supports of the fixed
old entries and the point such that the following holds.  Given a
list $Q$ of at most $\chi$ positions in $J\setminus V$ and at
most $\chi$ positions $R$ in $S$, one can move $Q$ to a list in
$S$, preserving all order comparisons and equalities relative to
$R$ and $T^*$, as well as its side of~$V$.

For completeness, the size requirement on this use of support
closure is bounded independently of~$\mu$.  If $\lambda<\mu$
bounds $\chi$, the initial supports, and $T^*$, the witnesses can
be chosen through $\chi^+$ stages inside a set of size at most
$2^\lambda$.  At each stage there are at most $2^\lambda$
realized requests over at most $\chi$ old positions, since
$(2^\lambda)^\chi=2^\lambda$.  Every at-most-$\chi$ subset of
the union occurs at an earlier stage.  Strong-limit closure keeps
the resulting set~$S$ below~$\mu$.

The subpresentation $E(S\cup V)$ is a strong submodel of $E(J)$
and has cardinality~$\mu$.  The complete list of elements generated
by $S$ is bounded, since the number of labels and $|S|$ are below
$\mu$.  Saturation for bounded $\K$-diagrams in
$E(S\cup V)$ therefore realizes the joint diagram of $z^*$ with
this entire list, producing~$z$.  The diagram on that parameter
list is the same in both presentations by naturalness.  Equalities
to its entries preserve the specified old assignments.

Choose representations of $z$ using positions $T\subseteq S\cup
V$.  Let $q$ be a test of at most $\chi$ parameters from $B$, with
its chosen supports $Q\subseteq J\setminus V$.  Apply support
closure with $R=T\cap S$; this set has size at most~$\chi$.
It gives a moved parameter list $q'$ whose terms use only~$S$.
The joint order pattern of the supports of $(z,q)$ agrees with
that of $(z,q')$: comparisons to $T\cap S$ are retained, and
comparisons to $T\cap V$ are forced by the retained side of the
convex strip.  The pattern of $(z^*,q)$ likewise agrees with that
of $(z^*,q')$, because comparisons with $T^*$ were retained.
Naturality and the diagram chosen in $E(S\cup V)$ give
\[
 D^{\K}(z,q)=D^{\K}(z,q')
             =D^{\K}(z^*,q')=D^{\K}(z^*,q).
\]
Here $q'$ need not belong to $B$: all terms on $S$ were included
in the saturation request.  This proves equality on tests of
size at most~$\chi$.  Both tuples have placed comparisons over
the original permitted layout, so its testing bound gives equality
on every bounded test, as asserted.
\end{proof}

\begin{proof}[Proof of Theorem~\ref{exact:geometric-boundary}]
Continue the induction on the number of axes.  The zero-axis and
inactive-top cases were addressed above.  Suppose the top is active
and exact lifting is available on fewer axes.  Refine the uniform
levels for Lemma~\ref{exact:joint-boundary}, the working packages,
and the indicated one-dimensional layout bounds.  The compatible
face maps have a bijective union $h:B^-\to B^+$ preserving all
bounded joint $\K$-diagrams in the tops.

Use Proposition~\ref{sys:tiling} and geometric categoricity to present
the two systems by tiled presentations.  Choose each tile and buffer
as $H_i\times\mathbb Q$, where $H_i$ is a well-order of cardinality
$\mu$.  In the top-only tile choose $H_i$ of order type
$\mu\cdot\mu$; its $\mu$ consecutive intervals of order type
$\mu$ give pairwise disjoint convex substrips of cardinality~$\mu$.
Every proper face is generated on the fixed point positions, all
buffers, and its designated tiles, so its chosen representations
omit the top-only tile.  The finite family of these designated parts
is exactly a parameter layout covered by
Lemma~\ref{exact:tiled-testing}.  That lemma also allows the bounded
old matched tuple as fixed data.  Use its testing bound for this
layout in the following placement and realization arguments.

Construct nested bounded tuple matches in alternating directions.
The invariant for a match $a^-\mapsto a^+$ is
\begin{equation}
 D^{\K}(a^-,b)=D^{\K}(a^+,h(b))
 \quad\text{for every bounded occurrence-tagged }b\text{ in }B^-.
 \label{exact:back-forth-invariant}
\end{equation}
It includes the equality comparisons needed for a partial
injection, and initially holds with the empty tuple by
Lemma~\ref{exact:joint-boundary}.

For a forward step, enlarge $a^-$ to a requested bounded tuple
$v^-$ in the source top.  Prescribe its joint comparisons with
$B^+$ by transport along~$h$, and prescribe the actual values
$a^+$ for the old coordinates.  Every bounded local request is
possible in the target top.  To see this, collect all its boundary
parameters into one bounded list.  The invariant matches that
list jointly with $a^-$.  Bounded diagram homogeneity, or the
equivalent saturation clause, then realizes the requested
extension $v^-$ over its target copy and the actual old tuple.

Apply Lemma~\ref{calc:compact-placement} to the chosen boundary
parameter set, keeping the old tuple fixed.  The underlying $E$-order
need not be well-ordered, so use its padded-order clause with the
coherent locally possible prescriptions on every bounded subset.
For each bounded desired test size, closure under the full
profile-count recursion in that lemma supplies a placement in an
enlargement of this $E$-order realizing all prescriptions through
that size.  Choose the test size at least the bound from
Lemma~\ref{exact:tiled-testing}, including the old tuple as fixed
data.  Any two such outputs agree on that bound, so the lemma
identifies their full placed comparisons.  Fix one output.  For any
bounded query choose another output at a test size including that
query; their equality proves that the fixed output also has the
prescribed value there.  Thus a bounded tuple $v^*$ has the required
placed comparison on every bounded query, with the previous
coordinate equalities included.  These are actual placements in
arbitrary enlargements of the $E$-order; no horizontal-term
representation of their new supports is assumed.

All boundary supports avoid the top-only tile, and the old matched
tuple has boundedly many supporting positions.  Some convex
substrip of cardinality~$\mu$ in that tile avoids those positions
as well.  Apply Lemma~\ref{exact:internal-realization} to this
substrip.  It realizes the placed comparison by a tuple $v^+$
inside the actual target, retaining $a^+$ and satisfying
\eqref{exact:back-forth-invariant} for the larger tuple.  This
completes the forward step.  Interchanging the two systems and
using $h^{-1}$ gives the backward step with the same invariant.

Since $\cf(\mu)=\aleph_0$, choose countable bounded exhaustions
of both tops and alternate these two steps so that both are
covered.  Include in the requested tuples the enumerations of
smaller strong submodels used in the working-package map
criterion.  The union map is onto and preserves every atomic
formula and its negation, because any finite tuple occurs in a
stage.  It is therefore a $\K$-isomorphism of the tops; the same
map criterion also supplies strongness directly.  Equality tests
in \eqref{exact:back-forth-invariant} show that its restriction
to each proper face is the prescribed~$h_s$.  It is consequently
an isomorphism of the whole geometric systems.

For a fixed finite shape the dimension induction used only its
finitely many proper restrictions.  The remaining refinements
were the uniform arity bounds for working diagrams, the fixed
tiling layouts, support closure, and the bounded tree costs.
These are the finite recipes recorded in
Section~\ref{close:section-uniform}, establishing the last
assertion of the theorem.
\end{proof}
